%% ===================================================================
%%  Prior — citation intent poster.  A0 PORTRAIT, 3 columns.
%%  Generated from POSTER_CONTENT.md.  Compile:  pdflatex poster.tex
%%  (Callum: you never need to edit this file — edit POSTER_CONTENT.md
%%   and say "rebuild the poster".)
%% ===================================================================
\documentclass[final]{beamer}
\usepackage[size=a0,orientation=portrait,scale=1.30]{beamerposter}
\usetheme{default}
\usefonttheme{professionalfonts}
\setbeamertemplate{navigation symbols}{}

\usepackage[T1]{fontenc}
\usepackage[utf8]{inputenc}
\usepackage{helvet}
\renewcommand{\familydefault}{\sfdefault}
\usepackage{tikz}
\usepackage{amsmath}
\usepackage{graphicx}
\usepackage[most]{tcolorbox}
\usetikzlibrary{positioning,arrows.meta,shapes.geometric,calc,fit,backgrounds}

%% ---------- palette (matches Prior's own viewers) ------------------
\definecolor{pBg}    {HTML}{FAF6EC}   % poster background, cream
\definecolor{pPanel} {HTML}{FFFFFF}   % block background
\definecolor{pInk}   {HTML}{2E3440}   % body text
\definecolor{pDim}   {HTML}{6B7686}   % secondary text
\definecolor{pRule}  {HTML}{DCD8CC}   % hairlines
\definecolor{pTeal}  {HTML}{0A9396}   % accent
\definecolor{pDeep}  {HTML}{005F73}   % headings
\definecolor{pRust}  {HTML}{BB3E03}   % contrast / alert
\definecolor{pSand}  {HTML}{EE9B00}   % highlight
\definecolor{pSlate} {HTML}{8896A6}   % background class

\setbeamercolor{background canvas}{bg=pBg}
\setbeamercolor{normal text}{fg=pInk}
\usebeamercolor[fg]{normal text}

%% ---------- shorthands ---------------------------------------------
\newcommand{\ttx}[1]{\texttt{#1}}
\newcommand{\clsbg}{\textcolor{pSlate}{\textbf{\texttt{background}}}}
\newcommand{\clsuse}{\textcolor{pTeal}{\textbf{\texttt{uses\_extends}}}}
\newcommand{\clscmp}{\textcolor{pRust}{\textbf{\texttt{compares\_contrasts}}}}
\newcommand{\num}[1]{\textcolor{pDeep}{\textbf{#1}}}

%% ---------- the block environment ----------------------------------
\newtcolorbox{pblock}[1]{
  enhanced, breakable=false,
  colback=pPanel, colframe=pRule, boxrule=1.2pt, arc=8pt,
  left=20pt, right=20pt, top=22pt, bottom=14pt,
  boxsep=0pt, before skip=0pt, after skip=18pt,
  title={\raggedright\bfseries\fontsize{33}{39}\selectfont #1},
  coltitle=white, colbacktitle=pDeep,
  fonttitle=\bfseries, toptitle=12pt, bottomtitle=12pt,
  attach boxed title to top left={xshift=-1.2pt, yshift=-1.2pt},
  boxed title style={colframe=pDeep, arc=8pt, boxrule=0pt,
                     left=22pt, right=22pt}
}

%% a lighter, untitled note box
\newtcolorbox{pnote}{
  enhanced, colback=pBg, colframe=pRule, boxrule=1pt, arc=6pt,
  left=16pt, right=16pt, top=11pt, bottom=11pt,
  before skip=8pt, after skip=8pt
}

%% ---------- figure-or-placeholder ----------------------------------
%% \figbox{path}{height}{hint shown while the file is missing}
\newcommand{\figbox}[3]{%
  \IfFileExists{#1}{%
    \begin{center}\includegraphics[width=\linewidth]{#1}\end{center}%
  }{%
    \begin{center}
    \begin{tikzpicture}
      \node[draw=pRule, line width=1.5pt, dashed, fill=pBg, rounded corners=6pt,
            minimum width=0.95\linewidth, minimum height=#2, align=center,
            text=pDim, font=\fontsize{22}{28}\selectfont]
        {\textbf{[ screenshot pending ]}\\[6pt]\texttt{\detokenize{#1}}\\[6pt]\itshape #3};
    \end{tikzpicture}
    \end{center}%
  }%
}
\newcommand{\figcap}[1]{\vspace{-2pt}{\fontsize{18}{22}\selectfont\color{pDim}\itshape #1}\vspace{4pt}}

%% ---------- body text size ------------------------------------------
\newcommand{\body}{\fontsize{21}{27}\selectfont}
\newcommand{\smallbody}{\fontsize{19}{24}\selectfont}

\newlength{\colw}
\setlength{\colw}{0.305\paperwidth}

\begin{document}
\body
\begin{frame}[t]{}
\vspace{-1.5cm}

%% ===================================================================
%%  HEADER
%% ===================================================================
\begin{center}
\begin{tcolorbox}[enhanced, colback=pDeep, colframe=pDeep, arc=10pt,
                  left=40pt,right=40pt,top=26pt,bottom=26pt, width=\textwidth]
\centering
{\fontsize{86}{96}\selectfont\bfseries\color{white} Why did this paper cite that one?}\\[14pt]
{\fontsize{46}{54}\selectfont\color{pSand} Typed citation edges for a research-literature knowledge graph}\\[22pt]
{\fontsize{30}{36}\selectfont\color{white}\bfseries Callum Young}\quad
{\fontsize{30}{36}\selectfont\color{pRule} $\cdot$\ \ \textit{affiliation to be filled in}}\\[16pt]
\textcolor{pSand}{\rule{0.5\textwidth}{2pt}}\\[14pt]
{\fontsize{30}{38}\selectfont\color{white}
 809 citation sites \ \ $\cdot$ \ \ 3-class intent taxonomy \ \ $\cdot$ \ \
 \textbf{95\% on a blind gold set} (macro-F1 0.892)}
\end{tcolorbox}
\end{center}

\vspace{22pt}

%% ===================================================================
%%  THREE COLUMNS
%% ===================================================================
\noindent
%% -------------------------------------------------------------------
%%  COLUMN 1
%% -------------------------------------------------------------------
\begin{minipage}[t]{\colw}\vspace{0pt}

\begin{pblock}{1 \ \ What Prior is, and why citations}
\textbf{Prior} reads a corpus of research papers and builds a map of them. It pulls out
each paper's \emph{contributions} --- the specific things that paper claims to have done
--- and draws typed, directed relations between the contributions of different papers, so
you can ask \emph{``what is actually known about X?''} and get an answer grounded in the
literature rather than a plausible-sounding summary.

\medskip
The map has a weak link. Those relations are drawn from \textbf{semantic similarity plus
an LLM's reading of two contribution statements}. The model never sees whether one paper
actually cited the other, or what it said when it did. The result is a map that agrees
with itself: \num{70\% of edges are \ttx{supports}}, and \ttx{contradicts} --- the label
you would most want a literature map for --- has been measured at about
\textcolor{pRust}{\textbf{36\% precision}}.

\medskip
Meanwhile the citation record sits unused. And a citation is a different kind of object:
$p_1$ citing $p_2$ is a \textbf{fact}; an LLM saying their ideas ``build on'' each other is
an \textbf{assertion}.

\medskip
But raw citations are blunt. In the graph, $p_1$ \emph{extending} $p_2$'s method, $p_1$
\emph{mentioning} $p_3$ in passing, and $p_1$ \emph{name-dropping} $p_4$ out of obligation
are all the same edge with the same weight. \textbf{The reason for a citation is the
missing variable.}

\begin{pnote}
\smallbody\centering
\textbf{So: type the citation edges $(p_1, p_2, t)$,\\ then use those facts to discipline
the assertions.}
\end{pnote}

\medskip
\figbox{figures/atlas.png}{9cm}{The Prior atlas --- zoomed out, topic clusters visible}
\figcap{Fig 1. Prior's semantic map of the 152-paper corpus. Relations here are drawn
without ever looking at a citation.}
\end{pblock}

\begin{pblock}{2 \ \ The taxonomy, and why ours}
We started from \textbf{RefWarden}, a citation-verification tool, which asks whether a
cited paper \ttx{exists}, whether it \ttx{supports} the claim it is attached to, and
whether it is \ttx{obligatory} or merely \ttx{helpful}. Good questions --- but they all ask
\emph{is this citation sound?} Prior needs to know \emph{why the citation is there at all.}

\medskip
So we added a primary \textbf{intent} axis, keeping support and priority as secondary
stamps:

\smallskip
\begin{itemize}\itemsep5pt
  \item[] \clsbg{} --- cited as context, prior art, or a passing mention.
  \item[] \clsuse{} --- the citing paper actually adopts or builds on the target.
  \item[] \clscmp{} --- the citing paper sets its own work \emph{specifically} against
    the target.
\end{itemize}

\begin{pnote}
\smallbody
\textbf{The selection rule:} a class only earns its place if the downstream system would
\emph{act differently} on it.
\end{pnote}

That single test does all the work. ACL-ARC's \ttx{uses} and \ttx{extends} collapse into
one class, because Prior treats them identically. Semantic Scholar's \ttx{result} is
dropped --- ambiguous, and it triggers no distinct action. \ttx{exists} is excluded because
every cited paper is already in the corpus, so it is trivially yes (parked for the
hallucinated-citation demo on AI-generated papers).

\medskip
\textbf{And why we could not just use an off-the-shelf label set:}
\textcolor{pRust}{\textbf{Semantic Scholar has no contrast class at all.}} Contrast is the
single most valuable signal for Prior, because \ttx{contradicts} is exactly the relation
the semantic graph gets wrong. The standard taxonomy cannot express the one thing we most
need.

\medskip
Intent is also not redundant with verification: \num{65\%} of \clscmp{} sites still
\ttx{support} their local claim --- so the support axis genuinely cannot see intent.

\bigskip
%% ---- FIG 2: taxonomy comparison -------------------------------------
\begin{center}
\begin{tikzpicture}[font=\fontsize{19}{23}\selectfont]
  \def\cw{5.9}
  \def\rh{2.1}
  \node[minimum width=\cw cm, minimum height=1.5cm, align=center, fill=pBg,
        text=pDim, draw=pRule, line width=1pt] at (0,0) {scite};
  \node[minimum width=\cw cm, minimum height=1.5cm, align=center, fill=pBg,
        text=pDim, draw=pRule, line width=1pt] at ({\cw},0) {ACL-ARC};
  \node[minimum width=\cw cm, minimum height=1.5cm, align=center, fill=pBg,
        text=pDim, draw=pRule, line width=1pt] at ({2*\cw},0) {Sem.\ Scholar};
  \node[minimum width=\cw cm, minimum height=1.5cm, align=center, fill=pDeep,
        text=white, draw=pDeep, line width=1pt] at ({3*\cw},0) {\textbf{Ours}};

  \foreach \i/\a/\b/\c in {
    1/{mentioning}/{background}/{background},
    2/{supporting}/{uses, extends}/{method},
    3/{contrasting}/{compare/\\contrast}/{\textcolor{pRust}{\textit{--- none ---}}},
    4/{---}/{motivation,\\future work}/{result} }
  {
    \node[minimum width=\cw cm, minimum height=\rh cm, align=center, fill=pPanel,
          draw=pRule, line width=1pt] at (0,{-\i*\rh}) {\a};
    \node[minimum width=\cw cm, minimum height=\rh cm, align=center, fill=pPanel,
          draw=pRule, line width=1pt] at ({\cw},{-\i*\rh}) {\b};
    \node[minimum width=\cw cm, minimum height=\rh cm, align=center, fill=pPanel,
          draw=pRule, line width=1pt] at ({2*\cw},{-\i*\rh}) {\c};
  }
  \node[minimum width=\cw cm, minimum height=\rh cm, align=center, fill=pBg,
        draw=pDeep, line width=1.6pt, text=pSlate] at ({3*\cw},{-\rh}) {\textbf{background}};
  \node[minimum width=\cw cm, minimum height=\rh cm, align=center, fill=pBg,
        draw=pDeep, line width=1.6pt, text=pTeal] at ({3*\cw},{-2*\rh}) {\textbf{uses\_}\\\textbf{extends}};
  \node[minimum width=\cw cm, minimum height=\rh cm, align=center, fill=pBg,
        draw=pDeep, line width=1.6pt, text=pRust] at ({3*\cw},{-3*\rh}) {\textbf{compares\_}\\\textbf{contrasts}};
  \node[minimum width=\cw cm, minimum height=\rh cm, align=center, fill=pBg,
        draw=pDeep, line width=1.6pt, text=pDim] at ({3*\cw},{-4*\rh}) {\textit{dropped}};

  \node[text=pRust, align=center, font=\fontsize{18}{22}\selectfont]
    at ({1.5*\cw},{-5*\rh-0.1})
    {the gap that forced our own judge $\nearrow$};
\end{tikzpicture}
\end{center}
\figcap{Fig 2. Four taxonomies. Classes merge when the downstream action is the same, and
drop when there is no distinct action. Semantic Scholar's missing contrast class is why an
off-the-shelf label set was not an option.}
\end{pblock}

\begin{pblock}{3 \ \ The pipeline}
The corpus is \num{152 papers} on AI-for-science, already carrying \num{581 contributions}
and \num{989 semantic edges} from Prior's existing pipeline.

\medskip
For each paper we extract the bibliography from its arXiv \LaTeX{} source, resolve every
reference back to a paper in the corpus, and pull the \textbf{passage around the citation},
marking the reference being typed as \ttx{[CITED:TARGET]} and its neighbours as
\ttx{[CITED]}. Intra-corpus citations only.

\medskip
That gives \num{525 citation edges across 809 claim-sites}, later extended to
\num{1,103 sites over 643 paper-pairs}. A \textbf{site} is one (edge, passage) pair; an
\textbf{edge} is one citing$\rightarrow$cited paper pair, rolled up from its sites.

\medskip
Everything ran on a Claude Code subscription rather than a metered API --- which is why
every runner is checkpointed, resumable, and stops cleanly when the usage window closes.

\bigskip
%% ---- FIG 3: pipeline strip ------------------------------------------
\begin{center}
\begin{tikzpicture}[font=\fontsize{18}{22}\selectfont,
   stage/.style={draw=pRule, line width=1.2pt, fill=pBg, rounded corners=5pt,
                 minimum width=3.9cm, minimum height=2.6cm, align=center, inner sep=3pt},
   fin/.style={draw=pDeep, line width=1.8pt, fill=pDeep, text=white, rounded corners=5pt,
               minimum width=4.4cm, minimum height=2.6cm, align=center, inner sep=3pt},
   ar/.style={-{Stealth[length=6mm,width=4mm]}, line width=1.5pt, draw=pTeal}]
  \node[stage] (a) at (0,0) {152\\papers};
  \node[stage, right=0.8cm of a] (b) {arXiv\\\texttt{.bbl}};
  \node[stage, right=0.8cm of b] (c) {resolve to\\corpus paper};
  \node[stage, right=0.8cm of c] (d) {citing\\passage};
  \node[fin,   right=0.8cm of d] (e) {typed edge\\$(p_1,p_2,t)$};
  \draw[ar] (a)--(b); \draw[ar] (b)--(c); \draw[ar] (c)--(d); \draw[ar] (d)--(e);
  \node[below=3mm of c, text=pDim, font=\fontsize{17}{20}\selectfont] {525 edges};
  \node[below=3mm of d, text=pDim, font=\fontsize{17}{20}\selectfont] {809 sites};
  \node[below=3mm of e, text=pDim, font=\fontsize{17}{20}\selectfont] {+ support, priority};
\end{tikzpicture}
\end{center}
\figcap{Fig 3. From LaTeX source to a typed edge. The judge sees the citing passage and the
cited paper's abstract --- nothing else.}
\end{pblock}

\end{minipage}
\hfill
%% -------------------------------------------------------------------
%%  COLUMN 2
%% -------------------------------------------------------------------
\begin{minipage}[t]{\colw}\vspace{0pt}

\begin{pblock}{4 \ \ Building the judge, and the prompt}
We reused RefWarden's judge \textbf{transport} --- the batching, the evidence assembly, the
JSON contract --- and swapped out the \textbf{rubric}. Support and priority came for free;
the intent axis was \emph{a prompt change, not new software}. The judge sees the cited
paper's abstract as evidence, plus the citing passage with the target marked.

\begin{pnote}
\smallbody
\textbf{The interesting part is that the constraint inverted.}\\[6pt]
RefWarden's hard problem is telling the model \textbf{what \emph{not} to do} --- don't
judge the neighbouring citations, don't use prior knowledge, abstain when the abstract
doesn't say.\\[6pt]
Intent has \textbf{no abstain option}: the model must commit to a class on every site.
Caution stops being safe --- it just becomes a systematic bias toward the majority class.
\end{pnote}

\medskip
Which is exactly what happened.

\medskip
%% ---- FIG 4: prompt version strip ------------------------------------
\begin{center}
\begin{tikzpicture}[font=\fontsize{19}{23}\selectfont,
  ver/.style={draw=pRule, line width=1.4pt, fill=pBg, rounded corners=6pt,
              text width=6.0cm, align=center, inner sep=8pt, minimum height=3.0cm},
  vfin/.style={draw=pDeep, line width=2pt, fill=pBg, rounded corners=6pt,
              text width=6.0cm, align=center, inner sep=8pt, minimum height=3.0cm},
  ar/.style={-{Stealth[length=6mm,width=4mm]}, line width=1.5pt, draw=pTeal}]

  \node[ver] (v2) at (0,0)
    {\textbf{v2}\\[3pt]{\fontsize{17}{21}\selectfont ``the target is a\\\emph{dependency}''}};
  \node[ver] (v3) at (7.8,0)
    {\textbf{v3}\\[3pt]{\fontsize{17}{21}\selectfont broaden \texttt{uses\_extends}\\+ baseline tie-breaker}};
  \node[vfin] (v4) at (15.6,0)
    {\textcolor{pDeep}{\textbf{v4 (final)}}\\[3pt]{\fontsize{17}{21}\selectfont grouping rule moved\\\emph{into} \texttt{background}}};

  \draw[ar] (v2)--(v3);
  \draw[ar] (v3)--(v4);

  \node[below=0.9cm of v2, text width=6.4cm, align=center, text=pRust,
        font=\fontsize{17}{21}\selectfont]
        {\textbf{fails:} too cautious,\\ underfires \texttt{uses\_extends}};
  \node[below=0.9cm of v3, text width=6.4cm, align=center, text=pRust,
        font=\fontsize{17}{21}\selectfont]
        {\textbf{fails:} overfires \texttt{compares\_}\\\texttt{contrasts} on \emph{group} critiques};
  \node[below=0.9cm of v4, text width=6.4cm, align=center, text=pTeal,
        font=\fontsize{17}{21}\selectfont]
        {\textbf{95\%} on the blind\\ gold set};
\end{tikzpicture}
\end{center}
\figcap{Fig 4. Three rubric versions, each fixing the previous one's failure mode.}

\medskip
\begin{itemize}\itemsep5pt
\item \textbf{v2 was too cautious and underfired} \clsuse{}. The rubric was asymmetric: it
  demanded the target be ``a dependency of the citing work'' for adoption, while accepting
  much weaker cues for \ttx{background}.
\item \textbf{v3 broadened it} --- the target ``provides an ingredient, method,
  formulation, metric, dataset or architectural block that the citing paper actively
  uses'' --- and added a baseline tie-breaker: \emph{rerunning} a system to beat it is
  \clscmp{}, \emph{adopting} its protocol or code is \clsuse{}.
\item \textbf{v3 then overfired} \clscmp{} \textbf{on group critiques.} \emph{``Another
  limitation is that current systems are limited to small-scale code experiments
  \ttx{[CITED:TARGET]}''} is a critique of a field, not of a paper.
\item \textbf{v4 fixed it by moving the rule, not adding one.} The ``do not over-use
  \ttx{compares\_contrasts}'' warning came \emph{out} of the negative clause and went
  \emph{into} the positive definition of \ttx{background}: if the target is grouped with
  other papers to describe a general limitation, it \textbf{is} background.
\end{itemize}

\begin{pnote}
\smallbody\centering
\textbf{Telling the model what a class \emph{is} beats telling it\\ what a neighbouring
class \emph{isn't}.}
\end{pnote}

\medskip
This was iterative hand-review rather than a controlled ablation --- but v4 is the version
that was then measured, blind, at 95\%.
\end{pblock}

\begin{pblock}{5 \ \ Evaluation}
Three validation stages, in increasing order of cost and of trust.

\medskip
\textbf{A. Semantic Scholar as a silver standard --- a negative result.} Free, no LLM. Only
237 of our 525 edges appear in S2's reference graph and \textcolor{pRust}{\textbf{only 10
carry an intent label at all}}: SciCite has not been run on 2025--26 papers. On the 9
comparable edges S2 agreed 6/9, and \emph{on all three disagreements S2 was the coarser
one}. This is why the next two stages were necessary.

\medskip
\textbf{B. A blind second judge.} A \emph{stronger} model (Opus~5) relabelled a stratified
100-site sample, given a \textbf{clean-room rubric} --- it never saw our tie-breakers, so
the two annotators' errors decorrelate --- and blind to our label.

\smallskip
%% ---- FIG 5: kappa stat strip ----------------------------------------
\begin{center}
\begin{tikzpicture}[font=\fontsize{20}{24}\selectfont,
  tile/.style={draw=pRule, line width=1.2pt, fill=pBg, rounded corners=6pt,
               minimum width=6.4cm, minimum height=3.6cm, align=center}]
  \node[tile] (k) at (0,0)
    {\fontsize{42}{48}\selectfont\color{pDeep}\textbf{0.79}\\[2pt]
     \fontsize{18}{22}\selectfont\color{pDim}Cohen's $\kappa$};
  \node[tile, right=0.6cm of k] (a)
    {\fontsize{42}{48}\selectfont\color{pDeep}\textbf{81\%}\\[2pt]
     \fontsize{18}{22}\selectfont\color{pDim}agreement,\\
     \color{pDim}reweighted};
  \node[tile, right=0.6cm of a] (d)
    {\fontsize{42}{48}\selectfont\color{pTeal}\textbf{14/14}\\[2pt]
     \fontsize{18}{22}\selectfont\color{pDim}disagreements on\\
     \color{pDim}known taxonomy\\
     \color{pDim}boundaries};
\end{tikzpicture}
\end{center}
\figcap{Fig 5. Second-judge cross-check --- 100 stratified sites, blind.}

\smallskip
The result that matters is not the number: \textbf{all 14 disagreements land on known
taxonomy boundaries, and none are random errors.} A judge that only argues on the
genuinely ambiguous cases is a credible judge.

\medskip
\textbf{C. A hand-labelled gold set --- the real number.} 122 sites labelled blind through
a purpose-built web app; the judge's verdicts and quality gates sat in the backing sheet
but were never sent to the browser. The queue is \textbf{blocked by sample type}:

\smallskip
\begin{itemize}\itemsep5pt
\item[--] an 80-site \ttx{random\_eval} block \emph{in random order}, so any labelled
  prefix is still a uniform sample and the headline accuracy is unbiased;
\item[--] a \ttx{disagreement} block of stage-B hard cases, \textbf{excluded from
  accuracy}, used only to referee the taxonomy forks;
\item[--] a stratified top-up to $\sim$27 per class for macro-F1.
\end{itemize}

\medskip
%% ---- FIG 6: accuracy panel -------------------------------------------
\begin{center}
\begin{tikzpicture}[font=\fontsize{20}{24}\selectfont]
  \node[align=center] at (0,0)
    {{\fontsize{72}{80}\selectfont\color{pDeep}\textbf{95.0\%}}\\[4pt]
     {\fontsize{21}{25}\selectfont\color{pDim}76 / 80 random-block sites}};
  \begin{scope}[shift={(-6.0,-3.8)}]
    \draw[pRule, line width=1.4pt] (0,0)--(12.0,0);
    \foreach \x/\l in {0/80, 3/85, 6/90, 9/95, 12/100}
      \draw[pRule, line width=1.4pt] (\x,-0.16)--(\x,0.16)
        node[below=5pt, text=pDim, font=\fontsize{17}{20}\selectfont] {\l\%};
    \draw[pTeal, line width=7pt] ({(87.8-80)/20*12},0.6)--({(98.0-80)/20*12},0.6);
    \fill[pDeep] ({(95.0-80)/20*12},0.6) circle (0.30);
    \node[text=pDim, font=\fontsize{17}{20}\selectfont] at (6.0,1.6) {95\% CI [87.8, 98.0]};
  \end{scope}
  \node[align=center, font=\fontsize{21}{26}\selectfont] at (0,-7.2)
    {\textbf{94.6\%} reweighted to the population class mix\\
     \textbf{macro-F1 = 0.892}};
\end{tikzpicture}
\end{center}

\smallskip
%% ---- FIG 6b: confusion matrix ----------------------------------------
\begin{center}
\begin{tikzpicture}[font=\fontsize{19}{23}\selectfont]
  \def\cw{4.6}\def\rh{2.1}
  \node[text=pDim, font=\fontsize{18}{22}\selectfont] at ({2*\cw},{1.75}) {\textbf{gold}};
  \node[text=pSlate, font=\fontsize{19}{23}\selectfont] at ({1*\cw},{0.85}) {\textbf{bg}};
  \node[text=pTeal,  font=\fontsize{19}{23}\selectfont] at ({2*\cw},{0.85}) {\textbf{use}};
  \node[text=pRust,  font=\fontsize{19}{23}\selectfont] at ({3*\cw},{0.85}) {\textbf{cmp}};
  \node[text=pSlate, anchor=east, font=\fontsize{19}{23}\selectfont] at ({0.44*\cw},{0}) {\textbf{bg}};
  \node[text=pTeal,  anchor=east, font=\fontsize{19}{23}\selectfont] at ({0.44*\cw},{-\rh}) {\textbf{use}};
  \node[text=pRust,  anchor=east, font=\fontsize{19}{23}\selectfont] at ({0.44*\cw},{-2*\rh}) {\textbf{cmp}};
  \node[text=pDim, rotate=90, font=\fontsize{18}{22}\selectfont] at ({-0.12*\cw},{-\rh}) {\textbf{judge}};
  \foreach \i/\j/\v/\s in {
     0/1/{78.3}/100, 0/2/{1.3}/8,  0/3/{2.6}/14,
     1/1/{0.3}/4,    1/2/{6.2}/32, 1/3/{0.3}/4,
     2/1/{1.0}/7,    2/2/{0.0}/2,  2/3/{10.0}/45 }
  {
    \node[minimum width=\cw cm, minimum height=\rh cm, align=center,
          fill=pTeal!\s, draw=pRule, line width=1pt]
      at ({\j*\cw},{-\i*\rh}) {\v\%};
  }
  \node[text=pDim, font=\fontsize{17}{20}\selectfont] at ({2*\cw},{-2.6*\rh})
    {population-scaled};
\end{tikzpicture}
\end{center}
\figcap{Fig 6. Headline accuracy with its interval, and the reweighted confusion matrix.
Sites on the same edge are not independent, so the true interval is a little wider than
binomial.}

\medskip
\textbf{And the referee result.} On the 12 hard cases where our judge and Opus disagreed,
gold sided with \textbf{ours 6, Opus 5, neither 1}. \emph{The stronger model was not the
more accurate one} --- and the split was clean: ours won every
\ttx{background}$\rightarrow$X call, Opus won every X$\rightarrow$\ttx{background} call.
\end{pblock}

\end{minipage}
\hfill
%% -------------------------------------------------------------------
%%  COLUMN 3
%% -------------------------------------------------------------------
\begin{minipage}[t]{\colw}\vspace{0pt}

\begin{pblock}{6 \ \ Folding citations into the semantic graph}
Prior now holds \textbf{two graphs} over the same 152 papers, answering different
questions. The citation graph is paper$\rightarrow$paper and it is a \textbf{fact}. The
semantic graph is contribution$\rightarrow$contribution and it is an \textbf{assertion}.
Rolling the 989 semantic edges up to paper-pairs and intersecting gives:

\smallskip
%% ---- FIG 7: overlap venn ---------------------------------------------
\begin{center}
\begin{tikzpicture}[font=\fontsize{21}{25}\selectfont]
  \begin{scope}[opacity=0.5, transparency group]
    \fill[pTeal]  (-1.9,0) circle (4.2);
    \fill[pSand]  ( 1.9,0) circle (3.7);
  \end{scope}
  \draw[pTeal, line width=2pt] (-1.9,0) circle (4.2);
  \draw[pSand!80!black, line width=2pt] ( 1.9,0) circle (3.7);
  \node[white, font=\fontsize{32}{36}\selectfont] at (-3.9,0.5) {\textbf{535}};
  \node[white, font=\fontsize{18}{22}\selectfont, align=center] at (-3.9,-0.8) {semantic\\only};
  \node[white, font=\fontsize{32}{36}\selectfont] at ( 3.8,0.5) {\textbf{406}};
  \node[white, font=\fontsize{18}{22}\selectfont, align=center] at ( 3.8,-0.8) {citation\\only};
  \node[text=pInk, font=\fontsize{32}{36}\selectfont] at (0,0.6) {\textbf{118}};
  \node[text=pInk, font=\fontsize{18}{22}\selectfont] at (0,-0.9) {both};
  \node[text=pTeal, font=\fontsize{22}{26}\selectfont] at (-3.9,4.9) {\textbf{semantic}};
  \node[text=pSand!70!black, font=\fontsize{22}{26}\selectfont] at (3.8,4.3) {\textbf{citation}};
\end{tikzpicture}
\end{center}
\figcap{Fig 7. Paper-pair overlap. The two graphs are complementary, not redundant.}

\smallskip
Citations alone miss parallel work that never cites; semantics alone invents relations.

\begin{pnote}
\smallbody
\textbf{The result that justifies the taxonomy:} of the 26 shared pairs typed \clscmp{},
\textcolor{pRust}{\textbf{only 4}} carry any semantic \ttx{contradicts} edge. The critique
signal is almost entirely invisible to the semantic pass --- so the contrast class is not a
refinement, it is \emph{new information}.
\end{pnote}

\medskip
%% ---- FIG 8: intent x relation crosstab -------------------------------
\begin{center}
\begin{tikzpicture}[font=\fontsize{19}{23}\selectfont]
  \def\cw{3.7}\def\rh{1.75}
  \node[text=pDim, align=center, font=\fontsize{17}{20}\selectfont] at ({1*\cw},{1.3}) {supports};
  \node[text=pDim, align=center, font=\fontsize{17}{20}\selectfont] at ({2*\cw},{1.3}) {builds\_on};
  \node[text=pDim, align=center, font=\fontsize{17}{20}\selectfont] at ({3*\cw},{1.3}) {refines};
  \node[text=pDim, align=center, font=\fontsize{17}{20}\selectfont] at ({4*\cw},{1.3}) {contradicts};
  \node[text=pSlate, anchor=east, font=\fontsize{18}{21}\selectfont] at ({0.45*\cw},{0}) {\textbf{background}};
  \node[text=pTeal,  anchor=east, font=\fontsize{18}{21}\selectfont] at ({0.45*\cw},{-\rh}) {\textbf{uses\_extends}};
  \node[text=pRust,  anchor=east, font=\fontsize{18}{21}\selectfont] at ({0.45*\cw},{-2*\rh}) {\textbf{compares\_contrasts}};
  \foreach \i/\j/\v/\s in {
    0/1/{57}/62, 0/2/{23}/26, 0/3/{2}/4,  0/4/{8}/10,
    1/1/{10}/16, 1/2/{9}/14,  1/3/{1}/3,  1/4/{0}/1,
    2/1/{19}/24, 2/2/{10}/14, 2/3/{0}/1,  2/4/{4}/8 }
  {
    \node[minimum width=\cw cm, minimum height=\rh cm, align=center,
          fill=pTeal!\s, draw=pRule, line width=1pt] at ({\j*\cw},{-\i*\rh}) {\v};
  }
  \draw[pRust, line width=3pt, rounded corners=3pt]
    ({4*\cw-\cw/2},{-2*\rh-\rh/2}) rectangle ({4*\cw+\cw/2},{-2*\rh+\rh/2});
  \node[text=pDim, font=\fontsize{17}{20}\selectfont, align=center]
    at ({2.5*\cw},{-2.9*\rh}) {shared pairs carrying $\geq$1 semantic edge of that relation};
\end{tikzpicture}
\end{center}
\figcap{Fig 8. Citation intent $\times$ semantic relation, on the 118 shared pairs.}

\medskip
\textbf{Direction is nearly free.} Citation direction agrees with the semantic arrow on
\num{66} pairs and disagrees on \num{52} --- the LLM's \ttx{builds\_on}/\ttx{refines} arrow
is close to a coin flip. On the audited \clsuse{} pairs, citation direction was correct
\num{33/33}.

\medskip
\textbf{We then read the divergent pairs by hand}, marking each as \emph{complementary},
\emph{semantic-wrong}, \emph{citation-wrong} or \emph{rollup-artifact}. The dominant
verdict was \textbf{complementary} --- the two signals describe different \emph{facets} of
a pair rather than disagreeing. The clearest case is MLEvolve\,$\rightarrow$\,AlphaEvolve,
genuinely \textbf{both} \ttx{builds\_on} \textbf{and} \ttx{contradicts}: it extends the
paradigm \emph{and} claims to beat it. A majority-vote rollup would have hidden that
tension entirely.

\begin{pnote}
\smallbody\centering
\textbf{Use citation signal as complementary evidence and as a\\ direction fix --- not as a
relabel.}\\[6pt]
Re-typed with citation evidence visible, \clsuse{} pairs became\\ \ttx{builds\_on}
\textbf{89\%} of the time --- without ever overwriting\\ the semantic layer.
\end{pnote}

\medskip
\figbox{figures/view_intent.png}{6cm}{Two-layer viewer --- both layers on, filtered to the overlap}
\figcap{Fig 9. Both graphs on one layout: solid arrows are citations coloured by intent,
dashed lines are the semantic relations.}
\end{pblock}

%% ---- SECTION 7 : comment out this whole block to drop the ideation work
\begin{pblock}{7 \ \ Graph vs.\ web ideation \textnormal{\textit{(in progress)}}}
\textbf{The question:} does an LLM exploring the Prior atlas propose better research
directions than the same LLM exploring the open web? Same seed topic, same task, same
output schema, same generator; only the environment differs. Judged blind and pairwise by a
stronger model, order-swapped.

\medskip
\textbf{Status: 10 ideas per arm generated; judging not yet built.} Three pilot findings are
already worth reporting:

\smallskip
\begin{enumerate}\itemsep5pt
\item \textbf{Blinding failed first.} 9 of 10 ideas leaked their arm --- the graph arm cited
  internal node IDs, the web arm said ``web search''. Fixed with an explicit blinding
  contract and an automated leakage gate before judging.
\item \textbf{The arms converged.} 3 of 4 seeds produced the \emph{same core idea} on both
  sides --- unsurprising in hindsight, since the corpus \textbf{is} the public literature
  the web arm retrieves.
\item \textbf{The cause: the graph arm wasn't using the graph.} Across four runs it called
  \ttx{get\_edges} once and the citation tools zero times --- using the atlas as a private
  search index over the same papers. The fix was not to \emph{force} traversal but to make
  structure \textbf{advertise itself}, surfacing typed-relation breakdowns and citation
  counts in every tool result. Traversal rose to 3--4 edge calls per run, and 3 of 4 seeds
  then diverged from the web arm.
\end{enumerate}

\smallskip
%% ---- FIG 10: arm comparison ------------------------------------------
\begin{center}
\begin{tikzpicture}[font=\fontsize{19}{23}\selectfont]
  \def\cw{4.6}\def\rh{1.9}
  \node[text=pDim, font=\fontsize{17}{20}\selectfont] at ({1*\cw},{1.25}) {explore calls};
  \node[text=pDim, font=\fontsize{17}{20}\selectfont] at ({2*\cw},{1.25}) {seconds};
  \node[text=pDim, font=\fontsize{17}{20}\selectfont] at ({3*\cw},{1.25}) {cost / idea};
  \node[text=pTeal, anchor=east, font=\fontsize{20}{24}\selectfont] at ({0.44*\cw},{0}) {\textbf{GRAPH}};
  \node[text=pSand!70!black, anchor=east, font=\fontsize{20}{24}\selectfont] at ({0.44*\cw},{-\rh}) {\textbf{WEB}};
  \foreach \i/\j/\v in {0/1/{13.6}, 0/2/{95}, 0/3/{\$0.46},
                        1/1/{3.9},  1/2/{111}, 1/3/{\$0.41}}
    \node[minimum width=\cw cm, minimum height=\rh cm, align=center,
          fill=pBg, draw=pRule, line width=1pt] at ({\j*\cw},{-\i*\rh}) {\v};
\end{tikzpicture}
\end{center}
\figcap{Fig 10. Comparable cost per idea --- the claim under test is \emph{better ideas at similar cost}.}
\end{pblock}
%% ---- end section 7 ---------------------------------------------------

\begin{pblock}{8 \ \ Takeaways}
\begin{itemize}\itemsep5pt
\item \textbf{A taxonomy chosen by downstream action, not by convention.} Three intent
  classes, each kept only because Prior would act differently on it --- validated blind at
  \num{95\% accuracy}, \num{macro-F1 0.892}.
\item \textbf{The rubric was the lever, not the model.} The constraint inverted from ``do
  not classify'' to ``must classify'', and the fix for over-firing was to define a class
  \emph{positively} rather than forbid its neighbour.
\item \textbf{The two graphs are complementary} --- 118 shared pairs of 1,059 --- and the
  contrast signal is nearly invisible to the semantic layer (4 of 26 pairs). Citation
  evidence belongs as \emph{evidence and direction}, not as a relabel.
\item \textbf{Next:} feed intent into the cartographer's labelling prompt and measure
  whether \ttx{contradicts} precision moves off 36\%.
\end{itemize}

\end{pblock}

\end{minipage}

\end{frame}
\end{document}
